\documentclass[preview]{standalone}

\usepackage{amsmath}
\usepackage{amssymb}
\usepackage{stellar}
\usepackage{definitions}

\begin{document}

\id{psicologia-sistema-nervoso-comportamento}
\genpage

\section{Sistema nervoso e comportamento}

\begin{snippetdefinition}{neuroscienza-definition}{Neuroscienza}
    La \emph{neuroscienza} è lo studio scientifico del cervello, del sistema
    nervoso e dei rapporti tra attività neurale, processi mentali e
    comportamento.
\end{snippetdefinition}

\begin{snippetdefinition}{neurone-definition}{Neurone}
    Il \emph{neurone} è una cellula specializzata nel ricevere, elaborare e
    trasmettere informazioni ad altre cellule del corpo.
\end{snippetdefinition}

\begin{snippetdefinition}{neuroni-sensoriali-definition}{Neuroni sensoriali}
    I \emph{neuroni sensoriali} sono neuroni che trasmettono informazioni dai
    recettori sensoriali verso il sistema nervoso centrale.
\end{snippetdefinition}

\begin{snippetdefinition}{neuroni-motori-definition}{Neuroni motori}
    I \emph{neuroni motori} sono neuroni che trasmettono comandi dal sistema
    nervoso centrale verso muscoli e ghiandole.
\end{snippetdefinition}

\begin{snippetdefinition}{interneuroni-definition}{Interneuroni}
    Gli \emph{interneuroni} sono neuroni che collegano tra loro altri neuroni,
    rendendo possibile l'elaborazione interna dei segnali nel sistema nervoso.
\end{snippetdefinition}

\begin{snippetexample}{riflesso-allontanamento-dolore-example}{Riflesso di allontanamento dal dolore}
    Nel riflesso di allontanamento dal dolore, i recettori del dolore vicini
    alla pelle vengono stimolati da un oggetto pungente. I neuroni sensoriali
    trasmettono il segnale al midollo spinale, dove un interneurone attiva i
    neuroni motori. Questi ultimi stimolano i muscoli della zona interessata,
    permettendo al corpo di allontanarsi dallo stimolo prima che la percezione
    cosciente del dolore sia pienamente elaborata.
\end{snippetexample}

\includesnpt[width=65\%,src=/snippet/static-psychology/pain-withdrawal-reflex.jpg]{centered-img}

\begin{snippetdefinition}{neurone-specchio-definition}{Neurone specchio}
    Un \emph{neurone specchio} è un neurone che si attiva sia quando un
    individuo esegue un'azione motoria sia quando osserva un altro individuo
    eseguire la stessa azione.
\end{snippetdefinition}

\begin{snippetdefinition}{input-eccitatorio-definition}{Input eccitatorio}
    Un \emph{input eccitatorio} è un segnale che aumenta la probabilità che un
    neurone produca una risposta.
\end{snippetdefinition}

\begin{snippetdefinition}{input-inibitorio-definition}{Input inibitorio}
    Un \emph{input inibitorio} è un segnale che riduce la probabilità che un
    neurone produca una risposta.
\end{snippetdefinition}

\begin{snippetdefinition}{potenziale-azione-definition}{Potenziale d'azione}
    Il \emph{potenziale d'azione} è un impulso nervoso che si attiva in un
    neurone, percorre l'assone e può portare al rilascio di neurotrasmettitori
    nella sinapsi.
\end{snippetdefinition}

\begin{snippetdefinition}{legge-tutto-o-nulla-definition}{Legge del tutto o del nulla}
    La \emph{legge del tutto o del nulla} è il principio secondo cui, una volta
    raggiunta la soglia di attivazione, il potenziale d'azione si produce con
    intensità completa; un aumento ulteriore della stimolazione non ne aumenta
    la dimensione.
\end{snippetdefinition}

\begin{snippetdefinition}{sinapsi-definition}{Sinapsi}
    La \emph{sinapsi} è lo spazio funzionale di comunicazione tra due neuroni,
    attraverso cui avviene la trasmissione delle informazioni.
\end{snippetdefinition}

\begin{snippetdefinition}{neurotrasmettitore-definition}{Neurotrasmettitore}
    Un \emph{neurotrasmettitore} è un messaggero chimico rilasciato da un
    neurone nello spazio sinaptico. Attraversando la sinapsi, può produrre un
    effetto eccitatorio o inibitorio sul neurone postsinaptico.
\end{snippetdefinition}

\begin{snippetexample}{catecolamine-example}{Catecolamine}
    Le \emph{catecolamine} sono una classe di sostanze chimiche che comprende
    neurotrasmettitori come la \emph{norepinefrina} e la \emph{dopamina}.
    Entrambe sono rilevanti nello studio di disturbi psicologici come i
    disturbi d'ansia e dell'umore.
\end{snippetexample}

\begin{snippetdefinition}{sistema-nervoso-centrale-definition}{Sistema nervoso centrale}
    Il \emph{sistema nervoso centrale}, o \emph{SNC}, è la parte del sistema
    nervoso formata dall'encefalo e dal midollo spinale. Integra e coordina le
    funzioni corporee, elabora i messaggi in entrata e invia comandi alle
    diverse parti del corpo.
\end{snippetdefinition}

\begin{snippetdefinition}{sistema-nervoso-periferico-definition}{Sistema nervoso periferico}
    Il \emph{sistema nervoso periferico}, o \emph{SNP}, è la parte del sistema
    nervoso formata dai neuroni che collegano il sistema nervoso centrale con
    gli organi, i muscoli e i recettori sensoriali del corpo.
\end{snippetdefinition}

\begin{snippet}{divisioni-sistema-nervoso}
    Il sistema nervoso periferico comprende il \emph{sistema nervoso somatico},
    che regola le azioni dei muscoli volontari, e il \emph{sistema nervoso
    autonomo}, che regola funzioni corporee non controllate in modo cosciente.
    Il sistema nervoso autonomo si divide a sua volta in \emph{sistema
    simpatico}, coinvolto nelle risposte alle situazioni di emergenza, e
    \emph{sistema parasimpatico}, coinvolto nelle funzioni ordinarie
    dell'organismo.
\end{snippet}

\includesnpt[width=60\%,src=/snippet/static-psychology/nervous-system-overview.jpg]{centered-img}

\end{document}
